\chapter{Conclusion}\label{cap:conclusion}

This work set out to test whether the circuit-wide sparsification that follows the activation of adult-born granule cells
\textit{in vivo}~\parencite{mchughAdultborn2022} can be accounted for by the DG-CA3 microcircuitry as it is currently
characterized. The question was posed because the existing computational literature on neurogenesis and pattern separation stops
at the DG~\parencite{chavlisDendrites2017, kimEffect2024, yangDynamic2025}, and so could not address effects reported downstream,
and because the candidate circuit mechanisms proposed for adult-born cells cannot be distinguished from one another in a model
that excludes CA3~\parencite{aimoneComputational2016}.

To address it, a spiking model of the DG-CA3 network was assembled from the Hippocampome.org knowledge
base~\parencite{wheelerHippocampomeorg2023} and extended with a population of iGCs whose parameters were fitted to published
electrophysiological recordings~\parencite{espositoNeuronal2005}. Two experiments were then run on it. In the first, the level of
neurogenesis and the EC$\to$iGC connection probability were varied while pattern separation and population sparsity were measured
in the DG and in CA3. In the second, the iGCs were activated directly and the response of every population was recorded, as in the
\textit{in vivo} experiment that motivated the work.

Three findings follow. The mGC population became sparser and separated its inputs better as neurogenesis and the EC$\to$iGC
connection probability increased, while the iGCs themselves never separated their own inputs and acted instead as integrators, so
that the granule cell population considered as a whole moved in the opposite direction to its mature component. The pCA3
population followed the mGCs rather than the iGCs, becoming both sparser and more separating, although it continued to converge
its inputs overall, as is expected of a region organized for autoassociation. Directly activating the iGCs recruited inhibitory
interneurons in both the DG and CA3 and raised population sparsity, reproducing the effect observed \textit{in vivo}. Decomposing
the gain in mGC separation showed that most of it came from the reduction in the number of active cells rather than from a change
in how correlated their outputs were, which identifies sparsification as the mechanism responsible rather than orthogonalization.

These results speak to the question of what adult-born granule cells are for. They are consistent with the view of the iGCs as
transient modulators of the surrounding circuit rather than as encoders in their own right~\parencite{folszAdultborn2023,
mchughAdultborn2022}, and they attribute the effect to the circuitry rather than to plasticity, since the model reproduces it with
synapses that never change. They also show that the effect is not confined to the region in which the new cells sit: a population
of only 5\% of the granule cells was enough to carry both the sparsification and the improvement in separation into CA3.

The conclusions are limited by the simplifications set out in Chapter~\ref{cap:discussion}. Neurons are reduced to single
compartments, so the dendritic processing known to control GC sparsity is absent~\parencite{chavlisDendrites2017}; the interneuron
repertoire comprises two types in the DG and one in CA3, against the much larger diversity described
anatomically~\parencite{degroInterneuron2022}; and the iGCs are represented as a single static population, so the sweep over their
EC$\to$iGC connection probability samples the maturation range rather than following cells through it. The most important of these
is the absence of synaptic plasticity, which is what prevents the model from addressing storage and recall in the recurrent CA3
network and therefore leaves the influence of adult-born cells on CA3 engrams open~\parencite{koSystems2025,
kopsickFormation2024}. Introducing plasticity, adding more interneuron types, and letting the iGC population mature over time are
the natural next steps for this work.

A population of iGCs making up 5\% of the GCs was enough to sparsen the mGC code, raise its pattern separation, and carry both
effects into the CA3, while separating none of its own inputs. Regulation of coding sparsity across both regions therefore rests
on the DG-CA3 microcircuitry alone, without synaptic plasticity, with the iGCs acting as modulators of activity rather than as
encoders.
